% ATTEMPT_STATUS: unresolved
% TOP_PROBLEM: {"problemId": "problem.separable-quotient-problem-for-banach-spaces", "canonicalTitle": "Separable Quotient Problem for Banach Spaces", "releaseRank": 405, "primaryDomain": "analysis_pde", "primaryDomainLabel": "Analysis & PDE", "displayStatus": "open_with_solved_subcases", "releaseStatus": "open_with_solved_subcases"}
% !TeX program = lualatex
\documentclass[11pt]{article}
\usepackage[margin=25mm]{geometry}
\usepackage{fontspec}
\setmainfont{Latin Modern Roman}
\usepackage{amsmath,amssymb,mathtools}
\usepackage{unicode-math}
\setmathfont{Latin Modern Math}
\usepackage{xurl,hyperref}
\hypersetup{colorlinks=true,urlcolor=blue}
\setcounter{secnumdepth}{0}
\setlength{\parindent}{0pt}
\setlength{\parskip}{5pt}
\setlength{\emergencystretch}{3em}
\begin{document}
\section*{\texorpdfstring{Separable Quotient Problem for Banach Spaces}{Separable Quotient Problem for Banach Spaces}}\label{405}
\subsection{Definitions and mathematical statement}
A Banach space is a complete normed real or complex vector space. For a closed linear subspace $Y\subset X$, the quotient $X/Y$ carries the complete quotient norm
\[
 \|x+Y\|_{X/Y}=\inf_{y\in Y}\|x-y\|_X.
\]
A space is separable if it has a countable dense subset. The separable quotient problem asks whether every infinite-dimensional Banach space $X$ has some closed $Y$ for which $X/Y$ is both infinite-dimensional and separable. Neither a finite-dimensional quotient nor a separable subspace of $X$ answers this quotient question. For an already separable infinite-dimensional $X$, the zero subspace is allowed; the universal statement includes nonseparable spaces without a chosen basis or additional geometric structure.
\par\smallskip\textit{Formulation and concept references:} \href{https://arxiv.org/abs/2604.11832}{[S1]}.

\subsection{Short English statement}
Does every infinite-dimensional Banach space admit an infinite-dimensional separable quotient by a closed subspace?
\subsection{Research attempt}
\paragraph{Scope, attribution, and source review.}
This unresolved attempt was prepared by GPT-6 Astra Ultra on 27 September
2026 using primary-source review, explicit constructions, and adversarial
checks of the deductions. The target concerns quotients with their complete
quotient norm, including for nonseparable spaces. Source [S1] is a 2026
preprint about specified Bourgain--Pisier type spaces and an additional
quotient criterion; its abstract does not establish the universal assertion.
No result claimed there is needed below. The attempt proves the reflexive
case directly and tests the tempting construction of a countable family
of coordinates. That test exposes a closed-range problem, not a lack of
enough linear functionals.

\paragraph{A complete subcase: reflexive spaces.}
Let $X$ be an infinite-dimensional reflexive Banach space, over either
$\mathbb R$ or $\mathbb C$. Choose linearly independent functionals
$f_1,f_2,\ldots\in X^*$ and put
\[
 Z=\overline{\operatorname{span}}\{f_j:j\geq1\}\subset X^*.
\]
The space $Z$ is separable, infinite-dimensional and reflexive. Define
$Q:X\to Z^*$ by $Qx(z)=z(x)$. It is bounded with norm at most one.
It is also onto: a functional $\phi\in Z^*$ extends by Hahn--Banach to
$\widetilde\phi\in X^{**}$ with the same norm, and reflexivity writes
$\widetilde\phi(z)=z(x)$ for some $x\in X$. Thus $Qx=\phi$.

Here $Z^*$ is separable. To check this sometimes omitted point, the unit
ball of $Z^*$ is weak-star compact and metrizable because $Z$ is separable.
Reflexivity makes its weak and weak-star topologies agree. Take a countable
weakly dense subset of that ball. Its rational convex combinations (using
real rational coefficients even in the complex case) are countable and
norm dense in the ball: otherwise Hahn--Banach separation would separate
a point from their closed convex hull, contradicting weak density.
Consequently $Z^*$ is norm separable.

The kernel
\[
 Y=Z^\perp=\{x:z(x)=0\text{ for every }z\in Z\}
\]
is closed. In fact the induced map $X/Y\to Z^*$ is an isometry, since
the norm-preserving extension above shows that every $\phi$ has a
representative $x$ with $\|x\|=\|\phi\|$, while the reverse inequality
follows from $\|Q\|\leq1$. This gives the required infinite-dimensional
separable quotient for every reflexive $X$. It includes nonseparable
Hilbert spaces. Reflexivity is used exactly to turn the Hahn--Banach
extension in $X^{**}$ into an element of $X$; deleting that step does
not prove the general problem.

\paragraph{Other elementary families.}
If $X$ has a complemented infinite-dimensional separable subspace $E$,
a bounded projection $P:X\to E$ gives $X/\ker P\cong E$. An arbitrary
separable subspace need not be complemented, so merely taking the
closed span of a countable independent sequence does not give this
argument. For $\ell^p(\Gamma)$, $1\leq p<\infty$, and for $c_0(\Gamma)$,
restriction to a countably infinite subset $\Gamma_0\subset\Gamma$
is onto the corresponding separable sequence space. Extension by zero
is a bounded right inverse. These assertions verify the quotient norm
and surjectivity, rather than just a coordinatewise map into a
separable ambient space.

\paragraph{Countably many coordinates always give a dense-range map.}
Every infinite-dimensional Banach space admits a biorthogonal sequence
$(x_j,f_j)$ with $f_j(x_i)=\delta_{ij}$. Here is an induction that needs
no basis for $X$. Having chosen finitely many pairs, their common
functional kernel is infinite-dimensional. Choose a nonzero vector
there outside the finite span of previous vectors, and use
Hahn--Banach on that finite-dimensional span plus the new vector to
produce the next functional vanishing on previous vectors. Normalize
the new pair to have value one. This maintains biorthogonality in both
directions.

Choose $a_j>0$ with $\sum_j a_j^2\|f_j\|^2<\infty$, for example
$a_j=2^{-j}/(1+\|f_j\|)$. Then
\[
 Tx=(a_jf_j(x))_{j\geq1}
\]
defines a bounded operator $T:X\to\ell^2$. The image contains every
finitely supported sequence, because $Tx_j=a_je_j$. Its range is
therefore dense and infinite-dimensional. Nevertheless $T$ is compact:
truncating at coordinate $N$ gives a finite-rank operator and
\[
 \|T-T_N\|\leq
 \left(\sum_{j>N}a_j^2\|f_j\|^2\right)^{1/2}\longrightarrow0.
\]
It cannot be onto $\ell^2$. If it were onto, the open mapping theorem
would put a fixed ball of $\ell^2$ inside the image of a bounded ball
of $X$, contradicting relative compactness and the noncompactness of
the infinite-dimensional unit ball.

This construction does yield an infinite-dimensional quotient
$X/\ker T$, but it does not show that this quotient is separable in
its quotient norm. The induced bijection onto $T(X)$ need not have
continuous inverse when $T(X)$ is given its inherited $\ell^2$ norm.
That distinction is exactly where the naive argument fails.

\paragraph{An explicit norm-topology countertest.}
Take $X=\ell^\infty$ and
\[
 Tx=(2^{-j}x_j)_{j\geq1}\in\ell^2.
\]
The map is bounded, compact and injective, and its range contains all
finite sequences. Thus the kernel is zero and the quotient is
$\ell^\infty$ itself, which is nonseparable. One elementary proof is
that all indicator sequences of subsets of $\mathbb N$ form an
uncountable set of pairwise distance one. The separable completion
of $T(X)$ in the weaker target norm has not become a separable
Banach quotient of the original space.

There is also a concrete failure of surjectivity: the sequence
$y_j=2^{-j}\sqrt j$ lies in $\ell^2$, while its only possible
preimage has coordinates $\sqrt j$ and is unbounded. This example
does not say that $\ell^\infty$ has no separable quotient. It says
that this particular countable-coordinate construction has not
produced one.

\paragraph{A quantitative criterion that would repair the construction.}
Suppose $(g_j)\subset X^*$ is bounded and weak-star null, so
$g_j(x)\to0$ for every $x$. Then
\[
 Qx=(g_j(x))_{j\geq1}
\]
is a bounded map into $c_0$. Its adjoint on $\ell^1=c_0^*$ is
$Q^*a=\sum_j a_jg_j$, with norm convergence. A sufficient condition
for surjectivity is
\begin{equation}\label{a405:lower}
 \left\|\sum_j a_jg_j\right\|\geq c\sum_j|a_j|
 \quad\text{for all finitely supported }a,\qquad c>0.
\end{equation}
The inequality extends to all $\ell^1$. To see why it gives onto,
Hahn--Banach separation shows that the norm closure of $Q(B_X)$
contains $cB_{c_0}$: its support function at $a\in\ell^1$ is
$\|Q^*a\|\geq c\|a\|_1$. For complex spaces apply the same real
separation argument with real parts of complex functionals.

Given $y\in c_0$, approximate it by $Qx_1$ with
$\|x_1\|\leq2\|y\|/c$ and error at most $\|y\|/2$.
Repeat for the residual with successive errors halved.
The resulting series $\sum x_j$ converges in $X$, and its image
is exactly $y$. This proves surjectivity and hence the desired
quotient when \eqref{a405:lower} is available.

Weak-star nullity and normalization alone do not imply that lower
bound. In $X=\ell^2$, coordinate functionals have norm one and are
weak-star null, but the norm of a sum of the first $N$ is
$\sqrt N$, rather than a fixed multiple of $N$. The resulting
map $\ell^2\to c_0$ is not onto. This does not contradict the
reflexive-case proof, whose quotient need not be $c_0$.

\paragraph{Verification, gap, and outcome.}
The proofs identify the topology of each image, the closedness of
each kernel and the exact use of completeness. The compact diagonal
map and the coordinate-functional test check two different failure
modes. No numerical test can replace the universal lower estimate
\eqref{a405:lower} or the reflexivity argument. The first remaining
gap is a mechanism giving an appropriate closed-range map into some
infinite-dimensional separable Banach space for an arbitrary
nonreflexive $X$. Requiring the particular target $c_0$ would be
stronger and is not imposed by the problem.
\textbf{UNRESOLVED.} The reflexive and complemented-subspace cases
are proved, and a proposed universal shortcut is disproved.
No general solution or novelty is claimed.

\subsection{Sources}
\begin{itemize}
\item[S1]Mazur's Separable Quotient Problem for Nonseparable Bourgain-Pisier spaces.
\url{https://arxiv.org/abs/2604.11832}.
\textit{Formulation source.}
\end{itemize}
\end{document}
